\documentclass[11pt, ukrainian]{article}
\usepackage[a4paper]{geometry}
\setlength\parindent{0mm}
\geometry{top=1.1cm,bottom=1.1cm,left=1.0cm,right=1.0cm}
\pagestyle{empty}
\usepackage[T2A]{fontenc}
\usepackage[utf8]{inputenc}
\usepackage[ukrainian]{babel}
\usepackage{graphicx}
\begin{document}
%begin
\begin{large}
\textbf{Test Alpha, варіант \No 1}\par

%beginitem
1. Відмітити все, що є однією правильною лексемою:
( 1 ) 8,3

( 2 ) <>

( 3 ) 8.3

( 4 ) FIFO

%enditem

%beginitem
2. Відмітити все, що є коректним оператором С++:
( 1 ) cout>>3;

( 2 ) int x+y=z;

( 3 ) cout<<3.4;

( 4 ) int f(int n);

%enditem

%beginitem
3. Що буде виведено за виконання фрагменту коду?

%code
cout << 6 * 6 * 6 * 6;
%code

%enditem
\end{large}
%end
\vspace{20mm}
%begin
\begin{large}
\textbf{Test Alpha, варіант \No 2}\par

%beginitem
1. Відмітити все, що є однією правильною лексемою:
( 1 ) [1]

( 2 ) =<

( 3 ) z13

( 4 ) >

%enditem

%beginitem
2. Відмітити все, що є коректним оператором С++:
( 1 ) int f(int a,b);

( 2 ) if a<b a=0;

( 3 ) x-=2;

( 4 ) int x='a'-'c';

%enditem

%beginitem
3. Що буде виведено за виконання фрагменту коду?

%code
cout << 5 % 5 % 5 % 5;
%code

%enditem
\end{large}
%end
\newpage
\end{document}

